% ============================================================================
% 第 1 章 绪论
% ============================================================================

\section{绪论}

伏羲是一套以 Rust 实现的本地优先 AI Agent 编排平台。它把用户的自然语言意图、顶层 Agent 玄女的对话决策、多个并发执行的门客（基于 Claude Code 与 Codex 等 headless agent CLI 启动的工作进程）以及它们之间的协作消息，统一视作可观测、可回放、可审计的事件序列；通过单一 SQLite WAL 真相源、Tokio broadcast 实时分发与 Google A2A 协议互通三套基础设施，把多 Agent 系统从「演示脚本」推到「日用平台」。本章梳理伏羲所处的技术语境与同类系统现状，明确论文要解决的问题与五条贡献，并给出全文章节安排。

\begin{figure}[H]
\centering
\includegraphics[width=0.96\textwidth]{drawio/fig-1-1-overview.pdf}
\caption{伏羲平台总体示意}
\label{fig:overview-hero}
\end{figure}

\subsection{研究背景及意义}

AI Agent 的现代理论基础源自 Wooldridge 提出的「自治、社会能力、反应性、主动性」四属性\parencite{wooldridge2009mas} 与 Russell 等的「感知—决策—行动」闭环\parencite{russell2020aima}。大模型成熟之后，这一抽象层次被迅速工程化：Wang 等的综述把现代 Agent 抽象为「规划、记忆、工具使用、行动反馈」四组件\parencite{wang2023llmagentsurvey}；ReAct\parencite{yao2022react}、Toolformer\parencite{schick2023toolformer}、MRKL\parencite{karpas2022mrkl} 让 Agent 具备显式工具调用能力；Voyager\parencite{wang2023voyager} 通过 skill library 沉淀跨任务经验；SWE-agent\parencite{yang2024sweagent} 与 OpenHands\parencite{wang2024openhands} 把执行环境与 Agent-Computer Interface 提升为一等公民。这些工作共同表明，工具接口、执行隔离与持续记忆必须作为系统组件设计，而不是对话循环的附属物。

工程使用中，单 Agent 难以同时承担需求澄清、计划拆解、代码实现与结果交付。CAMEL\parencite{li2023camel}、AutoGen\parencite{wu2023autogen}、MetaGPT\parencite{hong2023metagpt}、AutoGPT\parencite{significantgravitas2023autogpt} 等多 Agent 框架证明角色分工能提升任务可解释性，但也暴露三类共性工程问题：通讯协议不统一、任务生命周期与持久化缺位、运行时状态难追踪。对一个本地化运行、要被日用的多 Agent 平台而言，高性能分布式通讯不是可选锦上添花而是平台前提——Agent 执行过程会产生密集的工具调用与中间消息，需要低延迟通道及时反馈给用户；多 Agent 并发调度不能依赖轮询；Agent 输出的概率性需要事件日志、会话恢复与执行隔离来降低不可追溯风险。这与 Lamport 在 1978 年提出的事件因果序思想\parencite{lamport1978time} 在工程上同构：一个看不到事件历史的多 Agent 系统形同一个没有时钟的分布式系统。

2024 年以来 headless agent CLI 的兴起把上述需求具象化：\texttt{claude -p}、\texttt{codex exec} 等工具以「整段任务一次性输入、过程事件流式输出」的形态投入生产，可被 shell 脚本、CI 流水线或上层 Agent 调用。多 Agent 协作不再是同一进程内的对象互调，而是多个独立子进程之间的结构化消息调度，这与分布式系统「以消息为契约、以进程为隔离单元」的传统范式重合。本工作面向「LLM agent CLI 已成既定事实、缺一个把通讯底座当一等公民设计、能在普通笔记本上稳定承载日常工作流并可跨节点扩展的平台」这一需求展开。

为方便后文阅读，先约定术语：顶层调度 Agent 称「玄女」，对应 \texttt{fuxi-orchestrator}；执行任务的工作 Agent 称「门客」，适配器为 \texttt{fuxi-agent-cc}/\texttt{fuxi-agent-codex}；事件分发抽象称「事件总线」，对应 \texttt{EventBus}。Google 主导的 Agent2Agent 协议简称 A2A，Anthropic 的 Model Context Protocol 简称 MCP——前者描述 Agent 间协作，后者描述模型与工具的接口，二者面向不同语义层次。

\subsection{国内外研究现状}

围绕 AI Agent 编排已形成几条平行研究路线。\textbf{Agent 行为模式}以 ReAct\parencite{yao2022react}、Toolformer\parencite{schick2023toolformer}、MRKL\parencite{karpas2022mrkl} 为代表，关注单 Agent 推理—行动—观察闭环的能力上限，但通讯层退化为函数调用，不涉跨进程协作。\textbf{多 Agent 协作框架}以 AutoGen\parencite{wu2023autogen}、MetaGPT\parencite{hong2023metagpt}、CAMEL\parencite{li2023camel}、AutoGPT\parencite{significantgravitas2023autogpt} 为代表，证明角色分工有效，但绝大多数以 Python 框架或云端 SaaS 为形态：消息走进程内对象图、状态绑死生命周期、不能跨节点扩展。后续工作如 ChatDev\parencite{qian2024chatdev}、CrewAI\parencite{crewai_2024} 与 AutoGen v0.4\parencite{ag2_2024} 仍延续 Python+角色对话主流路线，与本文 Rust + 跨进程门客 + 跨节点 A2A 路线在底层选择上分道。

\textbf{Agent 协议与编排标准}是另一条线。MCP\parencite{anthropic2024mcp} 标准化模型与工具接口，A2A\parencite{a2aproject2025a2a} 用 \texttt{AgentCard}、\texttt{Task}、\texttt{Message} 等概念描述 Agent 间互操作语义。A2A 已有 Python、Go、Java、JavaScript、.NET 与 Rust（\texttt{a2a-rs}）六语言官方 SDK，Rust 生态内还有 \texttt{a2a-client}、\texttt{a2a-types} 等第三方 crate。本工作不依托「Rust 生态空白」立论，而是面向「与本地优先编排层和事件总线深度集成、并把 \texttt{input-required} 语义映射为平台级人工介入状态」这一具体场景实现一份与平台同源演进的 A2A 风格协议适配层。\textbf{分布式系统经典工作}——Lamport 逻辑时钟\parencite{lamport1978time}、Kafka 追加式日志\parencite{kreps2011kafka}、Raft\parencite{ongaro2014raft}、ZooKeeper\parencite{hunt2010zookeeper}——虽不针对 Agent，但其追加式日志、push 式分发与故障恢复思想在伏羲事件总线中被反复借鉴。Parrot\parencite{lin2024parrot} 在 LLM 推理服务层关心 GPU 集群效率，与本工作在多进程编排层的关注点正交但互补。

工程基础设施已经齐备：Tokio 提供异步运行时与 broadcast/mpsc channel\parencite{tokio2024broadcast}，Axum 基于 tower/hyper 提供 Web 框架\parencite{axum2024docs}，git worktree 提供轻量任务隔离\parencite{git2024worktree}，WebSocket\parencite{fette2011websocket} 与 SSE\parencite{whatwg2024sse} 支撑实时通讯。Rust 生态围绕 LLM 的 \texttt{llm-chain}、\texttt{rig}、\texttt{kalosm} 等库多停留在「单 Agent 调模型 + 工具调用」层级，没有覆盖多 Agent 协作、跨节点调度或事件溯源这类系统级议题。同类 Python 系统普遍存在三类共性盲区：GIL 与协程切换让多核难以充分利用、上下文绑死 Python 对象图崩溃即丢、没有跨节点能力。伏羲选择 Rust+多进程门客+跨节点 A2A 是有意识地避开这些盲区。

\subsection{研究目标与贡献}

本文研究目标是：在 Rust 工作区中实现一个本地优先、事件驱动、可观测、可跨节点扩展的多 AI Agent 协作平台，把通讯底座作为一等公民设计，并通过吞吐、延迟与压力测试给出可量化性能边界。围绕这一目标做出五项可独立审视的工程贡献。

\textbf{贡献一：与平台同源演进的 A2A 风格协议适配层。} \texttt{fuxi-a2a} crate 共 1321 行（核心 1063 + 测试 258），覆盖 \texttt{agent/getCard}、\texttt{tasks/send}、\texttt{tasks/sendSubscribe}、\texttt{tasks/get}、\texttt{tasks/cancel} 五条核心路径，沿用早期 A2A JSON-RPC binding。差异化贡献集中在两点：把规范已有的 \texttt{input-required} 中间态从协议字段提升为平台级人工介入信号——映射到 \texttt{Task::PendingApproval} 与 \texttt{ShelfStatus::AwaitingInput}；采用单端点 HTTP+SSE 设计避免客户端管理两条独立连接。本实现采用早期 binding，与当前官方 PascalCase 方法集和 \texttt{Part} one-of 语义不完全对齐，亦未与 \texttt{a2a-rs} 做完整 wire-level 互通测试，故不主张与 v1.0 规范完全对齐。

\textbf{贡献二：非阻塞事件总线与 lag 哨兵机制。} \texttt{fuxi-events} 把实时分发、持久化与历史回放统一在同一抽象内。\texttt{publish()} 路径采用「\texttt{try\_send} + 后台转交 + 哨兵告警」三步组合：广播路径用 Tokio broadcast 零拷贝扇出，持久化路径 \texttt{try\_send} 投递 mpsc 队列、拥塞时后台 spawn 任务承担等待；当 writer 待处理数超过 \texttt{lag\_threshold = 512} 时发出 \texttt{Custom \{ "event\_store\_lagged" \}} 哨兵事件，让运维侧把背压从「不可见的死锁」看成「可见的告警」。该设计兼顾「调用方永不阻塞」与「事件不丢」两个看似矛盾的约束，是对 Kafka push 模式\parencite{kreps2011kafka} 在本地优先场景下的工程化适配。

\textbf{贡献三：玄女—门客分层与不可绕过的编排层。} 系统在 Agent 角色上确立顶层调度（玄女）与执行 Agent（门客）的显式分工。玄女是用户对话面唯一入口，所有门客之间通讯必须经玄女知情，任何调度决策都发出对应事件让观测端透明。这条「编排层不可绕过」硬约束在 \texttt{Fuxi::shutdown\_agent} 等敏感操作中反复实施：所有 shutdown 路径必须比对玄女 \texttt{AgentId}，命中则 warn 并 noop，例外仅限显式命名 \texttt{shutdown\_xuannv\_for\_handoff} 的会话交接路径。Worker 选择函数沿「角色匹配、节点钉死、最低饱和度」三层过滤逐级缩小候选集。

\textbf{贡献四：三层沙箱 \texorpdfstring{$\times$}{x} WebSocket 反连门客执行模式。} 针对多 Agent 并发的两个核心难题——文件系统隔离与进程间双向通讯——本文提出 L1/L2/L3 三层沙箱模型（只读视图 / 任务级 worktree / 角色级持久沙箱），依赖 git worktree\parencite{git2024worktree} 提供互不干扰的工作树；同时提出 WS 反连模式——编排层在本机启 WebSocket 服务器，子进程通过 \texttt{--sdk-url} 反连回来，使编排层在任意时刻都能主动推送 \texttt{send\_message} 或 \texttt{cancel}，不必等门客 poll stdio。这两项设计共同把「多 Agent CLI 适配」收敛到统一框架，已支持 Claude Code 与 Codex 两套差异极大的 headless agent CLI。

\textbf{贡献五：跨节点扩展点的钩子化抽象与 sandbox 验证。} 跨节点协作所需的扩展点抽象为四个依赖反转 trait：\texttt{RecallSink}、\texttt{DistEnqueuer}、\texttt{NodeLoadProvider}、\texttt{ExtractorSpawner}——统一定义在低层 crate，由顶层 \texttt{fuxi-cli} 在 daemon 启动时注入，规避「memory 反向依赖 orchestrator」之类循环依赖。\texttt{Project} 数据结构新增 \texttt{host\_nodes: Vec<String>} 标记可执行节点，配合 \texttt{auto\_pin\_from\_project} 入口的双保险守卫保证用户显式 \texttt{@<node>} 不被自动逻辑覆盖；新增字段以 \texttt{\#[serde(default)]} 搭配反回归测试保持跨版本兼容。当前实现已在两节点 sandbox 与少量真实跨机部署中跑通端到端调度，更高节点数与不可靠链路下的故障转移留待后续工作（详见 §6.2）。

上述五项贡献围绕同一命题——把 Agent 间通讯的契约、事实流、调度边界与执行隔离同时做对——从协议层到平台层逐级展开。本文最终交付一份 13 crate、约 8.25 万行 Rust 代码的可执行工作区，配套 1363 个单元测试、27 个集成测试与持续运行于作者 ERP 项目的真实业务验证。

\subsection{论文组织结构}

全文按「问题—方案—实现—验证—回顾」展开，共五章。\textbf{第一章}（本章）阐述背景、研究现状、目标与五项贡献。\textbf{第二章 系统总体设计}从功能与非功能需求出发，给出整体分层架构与通讯模型形式化定义，明确事件驱动、单一真相源、不可绕过编排三条公理，并定义吞吐量、端到端延迟等性能指标。\textbf{第三章 系统设计与实现}沿 13 个 crate 模块依赖自下而上展开，覆盖事件总线、A2A 协议、玄女—门客调度、执行适配器与三层沙箱、跨节点通讯，把每个模块的关键数据结构、接口签名与代码片段一并呈现。\textbf{第四章 系统测试与结果分析}沿正确性、实时性、吞吐扩展性、参数敏感性、通讯层压力承受五个维度组织实验，吞吐采用 5-run median、延迟采用 500 样本统计 p50/p99。\textbf{第五章 总结与展望}回顾全文工作并讨论分布式一致性、记忆语义召回、Agent 适配器扩展、安全模型与用户界面五个方向的不足与后续路径。

\subsection{本章小结}

本章梳理 Agent 理论与 LLM 工程化两条线索，确立「把通讯底座作为一等公民」的核心立场；横向考察行为模式、协作框架、协议标准与分布式经典工作四条研究路线，指出现有同类系统在 Python 中心、单进程状态、跨节点缺失三方面的共性盲区；给出五项工程贡献——A2A 风格协议适配层、非阻塞事件总线与 lag 哨兵、玄女—门客分层、三层沙箱与 WS 反连、跨节点扩展点钩子化抽象。下一章把这些贡献映射到具体模块边界与形式化指标定义。
